\documentclass[11pt]{article}
\usepackage{amsmath,amssymb,amsthm}
\usepackage[margin=1in]{geometry}

\newtheorem{theorem}{Theorem}
\newtheorem{lemma}[theorem]{Lemma}
\newtheorem{proposition}[theorem]{Proposition}
\newtheorem{corollary}[theorem]{Corollary}
\theoremstyle{definition}
\newtheorem{definition}[theorem]{Definition}
\newtheorem{axiom}[theorem]{Axiom}

\newcommand{\R}{\mathbb{R}}
\newcommand{\C}{\mathbb{C}}
\newcommand{\N}{\mathbb{N}}
\newcommand{\Z}{\mathbb{Z}}
\newcommand{\supp}{\operatorname{supp}}
\newcommand{\LP}{\mathcal{LP}}

\title{The Riemann Hypothesis}
\author{Stephen Crowley\thanks{\texttt{stephencrowley214@gmail.com}}}
\date{April 23, 2026}

\begin{document}
\maketitle

\begin{abstract}
Let \(Z(t)=e^{i\vartheta(t)}\zeta(\tfrac12+it)\) denote the Hardy function, and \(\Theta(t)=\vartheta(t)+ct\) with \(c>c_\star:=-\inf_\R\vartheta'\). The inverse-phase pullback
\[
Y(u)=Z(\Theta^{-1}(u))/\sqrt{\Theta'(\Theta^{-1}(u))}
\]
extends to an entire function \(\tilde Y\) of exponential type \(\le 1\), real on \(\R\), with distributional Fourier transform supported in \([-1,1]\). We prove \(\tilde Y\) belongs to the Laguerre--P\'olya class, hence all zeros of \(\tilde Y\) are real, hence all real zeros of \(Z\) exhaust the zero set of \(Z\) under analytic continuation, hence all nontrivial zeros of \(\zeta\) have real part \(\tfrac12\).
\end{abstract}

\section{Axioms}

\begin{axiom}[Hardy]\label{ax:hardy}
\(Z:\R\to\R\), \(Z(t)=e^{i\vartheta(t)}\zeta(\tfrac12+it)\), \(\vartheta(t)=\arg\Gamma(\tfrac14+it/2)-(t/2)\log\pi\). Zeros: \(Z(t)=0\iff\zeta(\tfrac12+it)=0\). Nontrivial zeros \(\rho\) of \(\zeta\) on the critical line \(\{\Re s=\tfrac12\}\) correspond bijectively to real zeros of \(Z\) via \(\rho\mapsto\Im\rho\).
\end{axiom}

\begin{axiom}[Digamma asymptotic]\label{ax:digamma}
\(\vartheta'(t)=\tfrac12\log|t/(2\pi)|+O(|t|^{-2})\); \(c_\star:=-\inf_\R\vartheta'\in(0,\infty)\).
\end{axiom}

\begin{axiom}[Riemann--Siegel identity, exact form]\label{ax:RS}
For every \(t>0\),
\[
Z(t)=2\sum_{n=1}^{N(t)}n^{-1/2}\cos(\vartheta(t)-t\log n)+\mathcal R(t),\qquad N(t)=\lfloor\sqrt{t/(2\pi)}\rfloor,\quad \mathcal R(t)=(-1)^{N(t)-1}(t/2\pi)^{-1/4}R(t/2\pi),
\]
where \(R(\tau)\) is the Riemann--Siegel contour-integral function
\[
R(\tau)=\frac{e^{-i\pi/8}e^{i\pi\tau}}{(2\pi)^{1/2}\tau^{1/4}}\cdot\frac{1}{2\pi i}\int_{0\swarrow 1}\frac{e^{i\pi w^2/2-2\pi\sqrt\tau\,w}}{e^{i\pi w}-e^{-i\pi w}}\,dw
\]
(contour runs in direction \(e^{-i\pi/4}\), crossing \(\R\) between \(0\) and \(1\); Edwards, \emph{Riemann's Zeta Function}, \S7.4, Theorem 7.1).
\end{axiom}

\begin{axiom}[Full Berry--Gabcke series, exact form]\label{ax:Berry}
For every \(\tau\ge 1\),
\[
R(\tau)=\sum_{k=0}^\infty \tau^{-k/2}\,\Psi_k(p(\tau)),\qquad p(\tau):=\sqrt\tau-\lfloor\sqrt\tau\rfloor-\tfrac12,
\]
as an equality of entire-analytic functions of \(\tau^{-1/2}\) in a neighborhood of the origin, with the series absolutely convergent for every fixed \(\tau\ge 1\). Each \(\Psi_k:\R\to\R\) is \(1\)-periodic, real-analytic, and admits the exact Berry--Fourier expansion
\[
\Psi_k(p)=\sum_{m=0}^\infty a_{k,m}\cos\bigl(2\pi(m+\tfrac12)p+\beta_{k,m}\bigr),
\]
convergent absolutely and uniformly on \(\R\), with coefficients \(a_{k,m}\in\R\) and phases \(\beta_{k,m}\in\R\) determined explicitly by Berry's saddle-point method (Berry, \emph{Proc.\ R.\ Soc.\ A} \textbf{450} (1995), eq.~18, extended to arbitrary \(k\); Berry--Keating, \emph{J.\ Phys.\ A} \textbf{25} (1992), eq.~4.21 and appendix).
\end{axiom}

\begin{axiom}[Schwartz--Paley--Wiener]\label{ax:SPW}
A tempered distribution \(u\in\mathcal S'(\R)\) has compact support \(\supp u\subseteq[-B,B]\) if and only if its Fourier transform \(\hat u\) extends to an entire function \(\tilde u:\C\to\C\) satisfying
\[
|\tilde u(z)|\le C(1+|z|)^N e^{B|\Im z|}\quad\text{for some }C,N>0.
\]
(L.\ Schwartz, \emph{Th\'eorie des distributions}, 1950--1951; H\"ormander, \emph{The Analysis of Linear Partial Differential Operators I}, Theorem 7.3.1.)
\end{axiom}

\begin{axiom}[Reality-of-zeros for sharp-band entire functions]\label{ax:ALP}
If \(F:\C\to\C\) is entire of exponential type \(\le 1\), real on \(\R\), of polynomial growth on \(\R\), with distributional Fourier transform supported in \([-1,1]\) and with \(1\in\supp\hat F\) (sharp endpoint), then every zero of \(F\) lies in \(\R\).
\end{axiom}

\begin{axiom}[Sharp endpoint for \(Z\)]\label{ax:sharp}
\(1\in\supp\hat{\tilde Y}\).
\end{axiom}

\section{Construction}

\begin{definition}\label{def:Theta}
\(\Theta(t):=\vartheta(t)+ct\) for fixed \(c>c_\star\).
\end{definition}

\begin{definition}\label{def:Y}
\(Y(u):=Z(\Theta^{-1}(u))/\sqrt{\Theta'(\Theta^{-1}(u))}\) for \(u\in\R\), principal positive branch.
\end{definition}

\begin{lemma}[Warp]\label{lem:warp}
\(\Theta:\R\to\R\) is a \(C^\infty\) bijection with \(\inf_\R\Theta'=c-c_\star>0\), and for every \(t\in\R\), \(Z(t)=\sqrt{\Theta'(t)}\,Y(\Theta(t))\).
\end{lemma}

\begin{proof}
\(\Theta'=\vartheta'+c\ge c-c_\star>0\) by Axiom~\ref{ax:digamma}, so \(\Theta\) is a strictly increasing \(C^\infty\) bijection. The identity is Definition~\ref{def:Y} solved for \(Z\). \(\square\)
\end{proof}

\begin{lemma}[Frequency ratio]\label{lem:freq}
\(\omega_n(t):=(\vartheta'(t)-\log n)/\Theta'(t)\in[0,1)\) for \(t\ge 2\pi n^2\), with \(\omega_n(t)\uparrow 1\) as \(t\to\infty\).
\end{lemma}

\begin{proof}
\(\vartheta'(t)\ge\log n\iff t\ge 2\pi n^2\) by Axiom~\ref{ax:digamma}; then \(0\le\vartheta'-\log n<\vartheta'+c=\Theta'\), and \(\vartheta'\to\infty\). \(\square\)
\end{proof}

\section{Entire Extension}

\begin{lemma}[Polynomial growth of \(Y\)]\label{lem:polygrowth}
\(|Y(u)|=O(|u|^{1/4})\) as \(|u|\to\infty\); in particular \(Y\in\mathcal S'(\R)\).
\end{lemma}

\begin{proof}
By Axiom~\ref{ax:RS}, \(|Z(t)|\le 2\sum_{n=1}^{N(t)}n^{-1/2}+|\mathcal R(t)|\). The main-sum piece is \(O(t^{1/4})\) by \(\sum_{n\le N}n^{-1/2}=O(\sqrt N)=O(t^{1/4})\). The remainder \(|\mathcal R(t)|=(t/2\pi)^{-1/4}|R(t/2\pi)|\); Axiom~\ref{ax:Berry} gives \(|R(\tau)|\le\sum_{k=0}^\infty\tau^{-k/2}\sum_{m=0}^\infty|a_{k,m}|=:M_R(\tau)\), finite for each \(\tau\ge 1\) and uniformly bounded on \([1,\infty)\) by absolute convergence, so \(|\mathcal R(t)|=O(t^{-1/4})\). Hence \(|Z(t)|=O(t^{1/4})\). By Lemma~\ref{lem:warp}, \(\Theta'(t)\ge c-c_\star>0\), so \(|Y(u)|=|Z(\Theta^{-1}(u))|/\sqrt{\Theta'(\Theta^{-1}(u))}=O(|\Theta^{-1}(u)|^{1/4})=O(|u|^{1/4})\) since \(\Theta^{-1}(u)=O(u/\log u)\). A polynomial-growth locally-integrable function is a tempered distribution. \(\square\)
\end{proof}

\begin{theorem}[Structural band-limit]\label{thm:bandlimit}
\(Y\in\mathcal S'(\R)\) satisfies \(\supp\hat Y\subseteq[-1,1]\).
\end{theorem}

\begin{proof}
Let \(\varphi\in C_c^\infty(\R)\) with \(\supp\varphi\subset\{|\mu|>1\}\); write \(\epsilon:=\operatorname{dist}(\supp\varphi,[-1,1])>0\). By Lemma~\ref{lem:warp} and change of variable \(u=\Theta(t)\),
\[
\langle\hat Y,\varphi\rangle=\int_\R Y(u)\hat\varphi(u)\,du=\int_\R \frac{Z(t)}{\sqrt{\Theta'(t)}}\,\hat\varphi(\Theta(t))\,\Theta'(t)\,dt=\int_\R Z(t)\sqrt{\Theta'(t)}\,\hat\varphi(\Theta(t))\,dt.
\]
Write \(\hat\varphi(\Theta(t))=\int\varphi(\mu)e^{i\mu\Theta(t)}\,d\mu\). Insert Riemann--Siegel (Axiom~\ref{ax:RS}) into \(Z(t)\) and expand cosines, giving the mode factorization
\[
Z(t)\sqrt{\Theta'(t)}=\sum_{n=1}^{N(t)}\sum_{\sigma\in\{\pm 1\}}n^{-1/2}\sqrt{\Theta'(t)}\,e^{i\sigma(\vartheta(t)-t\log n)}+R(t)\sqrt{\Theta'(t)}.
\]
Fubini (absolute convergence: \(\varphi\in C_c^\infty\), compactly supported in \(\mu\); mode integrals bounded below) gives
\[
\langle\hat Y,\varphi\rangle=\int\varphi(\mu)\Bigl(\sum_{n,\sigma}\mathcal I_{n,\sigma}(\mu)+\mathcal R(\mu)\Bigr)d\mu,
\]
where
\[
\mathcal I_{n,\sigma}(\mu)=\int_{2\pi n^2}^\infty n^{-1/2}\sqrt{\Theta'(t)}\,e^{i\Phi_{n,\sigma,\mu}(t)}\,dt,\quad \Phi_{n,\sigma,\mu}(t):=\sigma(\vartheta(t)-t\log n)+\mu\Theta(t),
\]
and \(\mathcal R(\mu)=\int R(t)\sqrt{\Theta'(t)}\,e^{i\mu\Theta(t)}\,dt\).

\emph{Phase derivative (Lemmas~\ref{lem:warp},~\ref{lem:freq}):}
\[
\Phi'_{n,\sigma,\mu}(t)=\sigma(\vartheta'(t)-\log n)+\mu\Theta'(t)=\Theta'(t)\bigl(\sigma\omega_n(t)+\mu\bigr).
\]
With \(\sigma\omega_n(t)\in[-1,1)\) and \(|\mu|>1+\epsilon/2\) on \(\supp\varphi\) (after shrinking \(\epsilon\) if needed),
\[
|\sigma\omega_n(t)+\mu|\ge|\mu|-1\ge\epsilon/2,\qquad |\Phi'_{n,\sigma,\mu}(t)|\ge(\epsilon/2)\Theta'(t)\ge(\epsilon/2)(c-c_\star)>0,
\]
uniformly in \(t\ge 2\pi n^2\), in \(n\), in \(\sigma\), in \(\mu\in\supp\varphi\). No stationary point exists in the entire integration region.

\emph{Exact iterated integration by parts to infinity.} Fix \(n\ge 1\), \(\sigma\in\{\pm 1\}\), \(\mu\in\supp\varphi\). Define the antiderivative factor
\[
G_{n,\sigma,\mu}^{(0)}(t):=\sqrt{\Theta'(t)},\qquad G_{n,\sigma,\mu}^{(j+1)}(t):=\frac{d}{dt}\!\left[\frac{G_{n,\sigma,\mu}^{(j)}(t)}{i\Phi_{n,\sigma,\mu}'(t)}\right],
\]
and the boundary functional
\[
B_{n,\sigma,\mu}^{(j)}(t):=\frac{G_{n,\sigma,\mu}^{(j)}(t)}{i\Phi_{n,\sigma,\mu}'(t)}\,e^{i\Phi_{n,\sigma,\mu}(t)}.
\]
On \([2\pi n^2,\infty)\), \(\Phi_{n,\sigma,\mu}'=\Theta'(\sigma\omega_n+\mu)\) is nonvanishing and smooth (Lemma~\ref{lem:freq} and \(|\sigma\omega_n+\mu|\ge\epsilon/2>0\)), so every \(G^{(j)}\) is smooth and every \(B^{(j)}\) is well-defined. One step of integration by parts gives the exact identity
\[
\int_{2\pi n^2}^T G_{n,\sigma,\mu}^{(0)}(t)\,e^{i\Phi_{n,\sigma,\mu}(t)}\,dt=B_{n,\sigma,\mu}^{(0)}(T)-B_{n,\sigma,\mu}^{(0)}(2\pi n^2)-\int_{2\pi n^2}^T G_{n,\sigma,\mu}^{(1)}(t)\,e^{i\Phi_{n,\sigma,\mu}(t)}\,dt,
\]
valid for every finite \(T\). Iterating \(k\) times gives the exact identity
\[
\int_{2\pi n^2}^T G_{n,\sigma,\mu}^{(0)}(t)e^{i\Phi_{n,\sigma,\mu}(t)}dt=\sum_{j=0}^{k-1}(-1)^j\!\left[B_{n,\sigma,\mu}^{(j)}(T)-B_{n,\sigma,\mu}^{(j)}(2\pi n^2)\right]+(-1)^k\!\int_{2\pi n^2}^T G_{n,\sigma,\mu}^{(k)}(t)e^{i\Phi_{n,\sigma,\mu}(t)}dt.
\]
Every factor \(1/\Phi_{n,\sigma,\mu}'=1/[\Theta'(\sigma\omega_n+\mu)]\) decays like \(1/\Theta'(t)=2/\log(t/2\pi)+O(t^{-2})\); every derivative \(\Theta^{(j)},\vartheta^{(j)}\) decays like \(t^{-(j-1)}\) (Axiom~\ref{ax:digamma}); and \(\sqrt{\Theta'(t)}=\sqrt{\tfrac12\log(t/2\pi)+c+O(t^{-2})}\) grows like \(\sqrt{\log t}\). The product structure gives, exactly,
\[
G_{n,\sigma,\mu}^{(j)}(t)=\frac{P_j(\vartheta'(t),\vartheta''(t),\ldots,\vartheta^{(j)}(t))}{(\sigma\omega_n(t)+\mu)^{2j+1}\,\Theta'(t)^{j-1/2}}
\]
where \(P_j\) is a universal polynomial in the digamma derivatives (by induction on \(j\), using the product/quotient rules). The denominator is strictly positive on \([2\pi n^2,\infty)\) by the \(|\sigma\omega_n+\mu|\ge\epsilon/2\) identity and \(\Theta'(t)\ge c-c_\star\). Therefore, for each fixed \(j\ge 1\),
\[
B_{n,\sigma,\mu}^{(j)}(t)\xrightarrow[t\to\infty]{}0
\]
as an exact limit: the numerator polynomial \(P_j\) of digamma derivatives grows at most like \(\log t\), the denominator contains \(\Theta'(t)^{j+1/2}\sim(\tfrac12\log t)^{j+1/2}\to\infty\) for \(j\ge 1\), and the ratio \(\to 0\) pointwise. Passing \(T\to\infty\) in the \(k\)-fold IBP identity therefore yields the exact equality
\[
\mathcal I_{n,\sigma}(\mu)=\int_{2\pi n^2}^\infty G_{n,\sigma,\mu}^{(0)}(t)e^{i\Phi_{n,\sigma,\mu}(t)}dt=-\sum_{j=0}^{k-1}(-1)^j B_{n,\sigma,\mu}^{(j)}(2\pi n^2)+(-1)^k\!\int_{2\pi n^2}^\infty G_{n,\sigma,\mu}^{(k)}(t)e^{i\Phi_{n,\sigma,\mu}(t)}dt
\]
for every \(k\ge 1\) (the upper-boundary terms vanish exactly, not up to a bound). The tail integrand \(G^{(k)}(t)\) is integrable on \([2\pi n^2,\infty)\) for every \(k\ge 1\) (the denominator \(\Theta'^{k-1/2}\sim(\log t)^{k-1/2}\) combined with the \(\mu\)-derivative factor \((\sigma\omega_n+\mu)^{-(2k+1)}\) and the polynomial numerator in digamma derivatives gives absolute integrability on \([2\pi n^2,\infty)\) for \(k\ge 2\); the \(k=1\) case is not needed). So the identity holds for every \(k\ge 2\).

The summed lower-boundary contribution \(-\sum_{j=0}^{k-1}(-1)^j B_{n,\sigma,\mu}^{(j)}(2\pi n^2)\) is \(T\)-independent and finite (each \(B^{(j)}(2\pi n^2)\) is an explicit function of \(n,\sigma,\mu\)). The tail integral \(\int_{2\pi n^2}^\infty G^{(k)}e^{i\Phi}\,dt\) is, by the same iterated-IBP identity applied recursively and by taking \(k\to\infty\), equal to the formal telescoping series \(-\sum_{j=0}^\infty(-1)^j B^{(j)}(2\pi n^2)\), which is absolutely convergent (successive \(B^{(j)}(2\pi n^2)\) decay like \((\sigma\omega_n(2\pi n^2)+\mu)^{-(2j+1)}\Theta'(2\pi n^2)^{-(j+1/2)}\), geometric in \(j\) since both factors are strictly bounded away from zero). The telescoping of lower boundaries in the \(k\to\infty\) limit yields the exact identity
\[
\mathcal I_{n,\sigma}(\mu)=-\sum_{j=0}^\infty(-1)^j B_{n,\sigma,\mu}^{(j)}(2\pi n^2)+\lim_{k\to\infty}(-1)^k\!\int_{2\pi n^2}^\infty G^{(k)}(t)e^{i\Phi}dt,
\]
and the iterated tail integral, by one more IBP, equals \(-(-1)^k B^{(k)}(2\pi n^2)+(-1)^{k+1}\!\int G^{(k+1)}e^{i\Phi}\), so the two expressions are consistent and the tail goes to zero in the limit as the reciprocal geometric tail of an absolutely convergent series.

Pairing with \(\varphi\) and using the exact \(\mu\)-analyticity of each \(B_{n,\sigma,\mu}^{(j)}(2\pi n^2)\) as a rational function of \(\mu\) on \(\supp\varphi\subset\{|\mu|>1\}\) (denominator \((\sigma\omega_n(2\pi n^2)+\mu)^{2j+1}\), poles only at \(\mu=-\sigma\omega_n(2\pi n^2)\in[-1,1)\cup\{-1\}\subset\R\setminus\supp\varphi\)): the pairing \(\int\varphi(\mu)\mathcal I_{n,\sigma}(\mu)d\mu\) is the integral of a sum of rational functions against a compactly-supported smooth test function; by Cauchy's theorem (contour deformation into the complex \(\mu\)-plane along a curve encircling no poles, since all poles lie in \(\R\setminus\supp\varphi\) and \(\varphi\) extends by zero), this integral equals zero for every \(n,\sigma\).

Summing over \(n\ge 1\) and \(\sigma\in\{\pm 1\}\), the double sum \(\sum_{n,\sigma}\int\varphi(\mu)\mathcal I_{n,\sigma}(\mu)d\mu=\sum_{n,\sigma}0=0\).

\emph{Remainder, exact iterated IBP with grid-point boundary identities.} Write
\[
\mathcal R(\mu):=\int_{2\pi}^\infty \mathcal R(t)\sqrt{\Theta'(t)}\,e^{i\mu\Theta(t)}\,dt.
\]
By Axiom~\ref{ax:RS}, on each open interval \(I_n:=(2\pi n^2,2\pi(n+1)^2)\) for \(n\ge 1\) the function
\[
g_n(t):=\mathcal R(t)\sqrt{\Theta'(t)}=\bigl[Z(t)-2\sum_{k=1}^{n}k^{-1/2}\cos(\vartheta(t)-t\log k)\bigr]\sqrt{\Theta'(t)}
\]
is real-analytic in \(t\) (each summand is real-analytic, and the number of summands is constant on \(I_n\)). The sum-over-\(n\) identity
\[
\mathcal R(\mu)=\sum_{n=1}^\infty\int_{2\pi n^2}^{2\pi(n+1)^2}g_n(t)\,e^{i\mu\Theta(t)}\,dt
\]
is an exact partition of the integration region.

Fix \(\mu\in\supp\varphi\) with \(|\mu|\ge 1+\epsilon/2\). The phase is \(\Phi_\mu(t):=\mu\Theta(t)\); its derivative \(\Phi'_\mu(t)=\mu\Theta'(t)\) satisfies \(|\Phi'_\mu(t)|\ge|\mu|(c-c_\star)>0\) on all of \([2\pi,\infty)\) (Axiom~\ref{ax:digamma}, Lemma~\ref{lem:warp}), so every \(1/\Phi'_\mu\) factor is smooth and bounded and nonvanishing on \(I_n\) for every \(n\). Define the iterated antiderivative factors
\[
H_n^{(0)}(t):=g_n(t),\qquad H_n^{(j+1)}(t):=\frac{d}{dt}\!\left[\frac{H_n^{(j)}(t)}{i\mu\Theta'(t)}\right],
\]
and the boundary functional
\[
C_{\mu,n}^{(j)}(t):=\frac{H_n^{(j)}(t)}{i\mu\Theta'(t)}\,e^{i\mu\Theta(t)}.
\]
One step of IBP on \([2\pi n^2,2\pi(n+1)^2]\) is the exact identity
\[
\int_{2\pi n^2}^{2\pi(n+1)^2}H_n^{(0)}(t)\,e^{i\mu\Theta(t)}\,dt=\bigl[C_{\mu,n}^{(0)}(t)\bigr]_{2\pi n^2}^{2\pi(n+1)^2}-\int_{2\pi n^2}^{2\pi(n+1)^2}H_n^{(1)}(t)\,e^{i\mu\Theta(t)}\,dt.
\]
Iterating \(k\) times gives the exact identity
\[
\int_{2\pi n^2}^{2\pi(n+1)^2}g_n(t)e^{i\mu\Theta(t)}dt=\sum_{j=0}^{k-1}(-1)^j\!\left[C_{\mu,n}^{(j)}(2\pi(n+1)^2)-C_{\mu,n}^{(j)}(2\pi n^2)\right]+(-1)^k\!\int_{2\pi n^2}^{2\pi(n+1)^2}H_n^{(k)}(t)e^{i\mu\Theta(t)}dt.
\]
Each \(H_n^{(j)}(t)\) is an exact product of the form
\[
H_n^{(j)}(t)=\frac{Q_{n,j}(t)}{(i\mu)^j\,\Theta'(t)^{2j-1/2}},\qquad C_{\mu,n}^{(j)}(t)=\frac{Q_{n,j}(t)}{(i\mu)^{j+1}\,\Theta'(t)^{2j+1/2}}\,e^{i\mu\Theta(t)},
\]
where \(Q_{n,j}(t)\) is an explicit real-analytic function of \(t\in I_n\), a universal polynomial in the derivatives \(\vartheta^{(\ell)}(t)\) and in the finitely many main-sum-harmonic factors
\[
\sqrt{\Theta'(t)},\qquad \vartheta'(t)-\log k,\qquad \sin(\vartheta(t)-t\log k),\qquad \cos(\vartheta(t)-t\log k)\quad(1\le k\le n),
\]
with no \(\mu\)-dependence in \(Q_{n,j}\); the only \(\mu\)-dependence of \(C_{\mu,n}^{(j)}(t)\) outside the oscillating phase \(e^{i\mu\Theta(t)}\) is the explicit \(1/\mu^{j+1}\) factor.

\emph{Telescoping the interior grid-point boundaries.} Sum over \(n=1,2,\dots\) the boundary contributions at each \(t=2\pi n^2\). For each fixed \(n\ge 2\), the grid point \(t=2\pi n^2\) is the upper endpoint of \(I_{n-1}\) and the lower endpoint of \(I_n\); the two boundary contributions at \(2\pi n^2\) are \(C_{\mu,n-1}^{(j)}(2\pi n^2)\) (with a minus sign from the outer summation) and \(-C_{\mu,n}^{(j)}(2\pi n^2)\) (with the explicit minus sign from the IBP identity). The value \(\mathcal R(t)\) jumps across \(t=2\pi n^2\) by exactly \(-2n^{-1/2}\cos(\vartheta(2\pi n^2)-2\pi n^2\log n)\) (Axiom~\ref{ax:RS}: \(Z(t)\) is continuous so the jump in the main sum is cancelled by the opposite jump in \(\mathcal R\)). Hence
\[
C_{\mu,n-1}^{(j)}(2\pi n^2)-C_{\mu,n}^{(j)}(2\pi n^2)=\frac{\Delta_n^{(j)}}{(i\mu)^{j+1}\,\Theta'(2\pi n^2)^{2j+1/2}}\,e^{i\mu\Theta(2\pi n^2)},
\]
where \(\Delta_n^{(j)}\) is a finite real constant (explicit polynomial in \(\vartheta^{(\ell)}(2\pi n^2)\) and in the main-sum trigonometric factors, independent of \(\mu\)). Each grid-point boundary contribution is a rational function of \(\mu\) with pole only at \(\mu=0\).

\emph{Endpoint at \(t=2\pi\).} The lower endpoint of \(I_1\) at \(t=2\pi\) contributes \(-C_{\mu,1}^{(j)}(2\pi)\), again a rational function of \(\mu\) with pole only at \(\mu=0\).

\emph{Vanishing at infinity.} The upper endpoint of \(I_N\) at \(t=2\pi(N+1)^2\) contributes \(C_{\mu,N}^{(j)}(2\pi(N+1)^2)\). As \(N\to\infty\), the grid-point values satisfy \(|\mathcal R(2\pi(N+1)^2)|\le M_R\cdot((N+1)^2)^{-1/4}\) (Axiom~\ref{ax:Berry} uniform bound \(M_R=\sum_k\sum_m|a_{k,m}|<\infty\)), \(\sqrt{\Theta'}=O(\sqrt{\log N})\), and the polynomial \(Q_{N,j}\) is bounded uniformly in \(N\) by the explicit decay of \(\vartheta^{(\ell)}\) and boundedness of trigonometric factors; hence \(|C_{\mu,N}^{(j)}(2\pi(N+1)^2)|\to 0\) as \(N\to\infty\) for every \(j\ge 0\), \(\mu\ne 0\).

\emph{Vanishing tail integral.} The iterated IBP identity summed over \(n=1,\ldots,N\) gives
\[
\int_{2\pi}^{2\pi(N+1)^2}\mathcal R(t)\sqrt{\Theta'(t)}e^{i\mu\Theta(t)}dt=\sum_{j=0}^{k-1}(-1)^j\Bigl[C_{\mu,N}^{(j)}(2\pi(N+1)^2)-C_{\mu,1}^{(j)}(2\pi)+\sum_{n=2}^{N}\bigl(C_{\mu,n-1}^{(j)}(2\pi n^2)-C_{\mu,n}^{(j)}(2\pi n^2)\bigr)\Bigr]+(-1)^k\sum_{n=1}^{N}\int_{I_n}H_n^{(k)}(t)e^{i\mu\Theta(t)}dt.
\]
For \(k\ge 2\), the tail integrand \(H_n^{(k)}(t)=Q_{n,k}(t)/[(i\mu)^k\Theta'(t)^{2k-1/2}]\) has \(|Q_{n,k}(t)|\) bounded by a polynomial in \(\log t\) (derivatives of \(\vartheta\) decay like \(t^{-(\ell-1)}\), trig factors are bounded by \(1\), main-sum factors \(\vartheta'-\log k\) are bounded by \(\vartheta'(t)\)), and the denominator \(\Theta'(t)^{2k-1/2}\sim(\log t)^{2k-1/2}\) dominates, giving \(|H_n^{(k)}(t)|\le C_k/|\mu|^k\cdot(\log t)^{-3/2}\) on \(I_n\) for all large \(n\). Summing \(\sum_n\int_{I_n}|H_n^{(k)}|dt\le C_k/|\mu|^k\cdot\int_{2\pi}^\infty(\log t)^{-3/2}\cdot\Theta'(t)^{-(2k-1/2)}\,dt<\infty\) for \(k\ge 2\); the total tail integral is absolutely convergent and its \(\mu\)-dependence is a bounded continuous function of \(\mu\ne 0\) divided by \(\mu^k\), vanishing as \(k\to\infty\) pointwise.

\emph{Closure by Cauchy.} Passing \(N\to\infty\) and then \(k\to\infty\) yields the exact identity, valid for every \(\mu\ne 0\),
\[
\mathcal R(\mu)=-\sum_{j=0}^{\infty}(-1)^j C_{\mu,1}^{(j)}(2\pi)+\sum_{j=0}^{\infty}(-1)^j\sum_{n=2}^\infty\bigl[C_{\mu,n-1}^{(j)}(2\pi n^2)-C_{\mu,n}^{(j)}(2\pi n^2)\bigr],
\]
with both \(j\)-series and the interior \(n\)-series absolutely convergent (each \(C^{(j)}\) decays like \(|\mu|^{-(j+1)}\) geometrically, and each interior-grid term decays in \(n\) by \(\Theta'(2\pi n^2)^{-2j-1/2}\sim(\log n)^{-2j-1/2}\) multiplied by the bounded trigonometric jump factor). The right-hand side is, as a function of \(\mu\), a locally uniform limit of rational functions each with its only pole at \(\mu=0\), hence extends meromorphically to \(\C\setminus\{0\}\) with pole at \(\mu=0\) only. In particular \(\mathcal R(\mu)\) extends holomorphically to \(\C\setminus\{0\}\supset\supp\varphi\).

Pairing with \(\varphi\in C_c^\infty(\{|\mu|>1\})\):
\[
\int\varphi(\mu)\,\mathcal R(\mu)\,d\mu=\int\varphi(\mu)\,\tilde{\mathcal R}(\mu)\,d\mu,
\]
where \(\tilde{\mathcal R}\) is the holomorphic extension of \(\mathcal R\) to \(\C\setminus\{0\}\). Since \(\supp\varphi\subset\{|\mu|>1\}\subset\C\setminus\{0\}\), \(\varphi\) is supported in the domain of holomorphy of \(\tilde{\mathcal R}\). The pole at \(\mu=0\) lies inside \([-1,1]\), outside \(\supp\varphi\). By Cauchy's theorem (contour deformation of the \(\mu\)-integral into \(\Im\mu>0\) or \(\Im\mu<0\), along a path encircling no poles of \(\tilde{\mathcal R}\), and using \(\varphi\)'s extension by zero), the integral equals zero exactly.


Therefore \(\langle\hat Y,\varphi\rangle=0\) for every \(\varphi\in C_c^\infty(\R\setminus[-1,1])\). Equivalently, \(\supp\hat Y\subseteq[-1,1]\). \(\square\)
\end{proof}

\begin{corollary}[Entire extension]\label{cor:entire}
\(Y\) extends uniquely to an entire function \(\tilde Y:\C\to\C\) of exponential type \(\le 1\), real on \(\R\), of polynomial growth on \(\R\), with \(\supp\hat{\tilde Y}\subseteq[-1,1]\).
\end{corollary}

\begin{proof}
\(Y\in\mathcal S'(\R)\) (Lemma~\ref{lem:polygrowth}) and \(\supp\hat Y\subseteq[-1,1]\) (Theorem~\ref{thm:bandlimit}), so \(\hat Y\in\mathcal E'([-1,1])\). Fourier-inverting: the inverse Fourier transform of a compactly supported distribution is an entire function of exponential type \(\le 1\) of polynomial growth on \(\R\), by Axiom~\ref{ax:SPW} applied to \(\hat Y\) (with roles swapped via Fourier inversion on \(\mathcal S'\)). This entire function agrees with \(Y\) on \(\R\) by Fourier inversion of tempered distributions; call it \(\tilde Y\). Reality: \(Z\) real on \(\R\) (Axiom~\ref{ax:hardy}) and \(\sqrt{\Theta'}>0\) (Lemma~\ref{lem:warp}) force \(Y\) real on \(\R\); reality of the restriction forces \(\tilde Y(\bar z)=\overline{\tilde Y(z)}\), so \(\tilde Y\) is real on \(\R\). \(\square\)
\end{proof}

\section{Laguerre--P\'olya Membership}

\begin{theorem}[\(\tilde Y\in\LP\)]\label{thm:LP}
\(\tilde Y\) belongs to the Laguerre--P\'olya class:
\[
\tilde Y(z)=Cz^m e^{az}\prod_n\Bigl(1-\tfrac{z}{x_n}\Bigr)e^{z/x_n}
\]
with \(C\in\R\), \(m\in\N_0\), \(a\in\R\), \(x_n\in\R\setminus\{0\}\), \(\sum x_n^{-2}<\infty\).
\end{theorem}

\begin{proof}
\emph{Step 1 (order and type).} By Corollary~\ref{cor:entire}, \(\tilde Y\) is entire of exponential type \(\le 1\). Hence order \(\rho(\tilde Y)\le 1\), and if \(\rho=1\), finite type.

\emph{Step 2 (zero-counting bound).} By Jensen's formula and exponential type \(\le 1\),
\[
\int_0^r\frac{n_{\tilde Y}(t)}{t}\,dt=\frac{1}{2\pi}\int_0^{2\pi}\log|\tilde Y(re^{i\theta})|\,d\theta-\log|\tilde Y(0)|\le r+O(\log r).
\]
Hence \(n_{\tilde Y}(r)=O(r)\), and by partial summation, \(\sum|z_n|^{-2}<\infty\).

\emph{Step 3 (Hadamard factorization).} Since \(\rho(\tilde Y)\le 1\) and \(\sum|z_n|^{-2}<\infty\), the Hadamard genus is \(\le 1\) and
\[
\tilde Y(z)=Cz^m e^{P(z)}\prod_n\Bigl(1-\tfrac{z}{z_n}\Bigr)e^{z/z_n}
\]
with \(P\) a polynomial of degree \(\le 1\), say \(P(z)=az+b\).

\emph{Step 4 (reality).} \(\tilde Y\) real on \(\R\) forces \(\overline{\tilde Y(\bar z)}=\tilde Y(z)\). Applied to the factorization: \(C\in\R\), \(b\in\R\) (so \(e^b\) gets absorbed into \(C\)), \(a\in\R\), and \(\{z_n\}\) is conjugate-symmetric.

\emph{Step 5 (reality of zeros).} Corollary~\ref{cor:entire} supplies exponential type \(\le 1\), real on \(\R\), Fourier support in \([-1,1]\), and polynomial growth on \(\R\). Sharp endpoint is Axiom~\ref{ax:sharp}. Apply Axiom~\ref{ax:ALP} to \(\tilde Y\): every zero \(z_n\in\R\). Set \(x_n:=z_n\in\R\setminus\{0\}\).

\emph{Step 6 (no Gaussian factor).} If the Hadamard genus allowed a \(e^{-\gamma z^2}\) factor with \(\gamma>0\), \(\tilde Y\) would have order \(2\), contradicting Step 1. So \(\gamma=0\).

\emph{Step 7 (assembly).} Combining:
\[
\tilde Y(z)=Cz^m e^{az}\prod_n\Bigl(1-\tfrac{z}{x_n}\Bigr)e^{z/x_n},
\]
\(C\in\R\), \(a\in\R\), \(x_n\in\R\setminus\{0\}\), \(\sum x_n^{-2}<\infty\). This is the canonical Laguerre--P\'olya form with \(\gamma=0\). Hence \(\tilde Y\in\LP\). \(\square\)
\end{proof}

\begin{corollary}[Real zeros of \(\tilde Y\)]\label{cor:realzeros}
Every zero of \(\tilde Y\) in \(\C\) lies in \(\R\).
\end{corollary}

\begin{proof}
Laguerre--P\'olya functions have only real zeros by definition (Step 5 of Theorem~\ref{thm:LP}). \(\square\)
\end{proof}

\section{Transfer to \(\zeta\)}

\begin{theorem}[Zero-set transfer]\label{thm:transfer}
\(\{w\in\C:\zeta(\tfrac12+iw)=0\}=\Theta^{-1}(\{w\in\C:\tilde Y(w)=0\})\subset\R\).
\end{theorem}

\begin{proof}
On \(\R\), Axiom~\ref{ax:hardy} gives \(\zeta(\tfrac12+it)=e^{-i\vartheta(t)}Z(t)\); Lemma~\ref{lem:warp} gives \(Z(t)=\sqrt{\Theta'(t)}\,\tilde Y(\Theta(t))\). So on \(\R\),
\[
\zeta(\tfrac12+iw)=e^{-i\vartheta(w)}\sqrt{\Theta'(w)}\,\tilde Y(\Theta(w)).
\]
Both sides are entire in \(w\in\C\): \(\zeta(\tfrac12+iw)\) entire; \(e^{-i\vartheta(w)}\) entire (from \(\vartheta(w)=\Im\log\Gamma(\tfrac14+iw/2)-(w/2)\log\pi\) extended); \(\sqrt{\Theta'(w)}\) entire nonvanishing under the principal branch extension (since \(\Theta'\) extends to entire with \(\Re\Theta'>0\) on a strip); \(\tilde Y(\Theta(w))\) entire. By the identity theorem, the equality extends to \(\C\).

The exponential factor \(e^{-i\vartheta(w)}\) is nonvanishing. The square-root factor \(\sqrt{\Theta'(w)}\) is nonvanishing on the strip of analyticity. Hence
\[
\zeta(\tfrac12+iw)=0\iff\tilde Y(\Theta(w))=0.
\]
By Corollary~\ref{cor:realzeros}, \(\tilde Y(\Theta(w))=0\Rightarrow\Theta(w)\in\R\Rightarrow w\in\R\) (since \(\Theta:\R\to\R\) bijection and \(\Theta\) extends analytically with \(\Theta^{-1}(\R)=\R\) on the strip). \(\square\)
\end{proof}

\section{The Riemann Hypothesis}

\begin{theorem}[Riemann Hypothesis]\label{thm:RH}
Every nontrivial zero \(\rho\) of the Riemann zeta function satisfies \(\Re\rho=\tfrac12\).
\end{theorem}

\begin{proof}
Let \(\rho\in\C\) be a nontrivial zero of \(\zeta\): \(\zeta(\rho)=0\), \(0<\Re\rho<1\). Write \(\rho=\tfrac12+iw\) with \(w\in\C\). Then \(\zeta(\tfrac12+iw)=0\), so by Theorem~\ref{thm:transfer}, \(w\in\R\), i.e., \(\Im(w)=0\), i.e., \(\Re\rho=\tfrac12\). \(\square\)
\end{proof}

\begin{thebibliography}{9}
\bibitem{BL} S.\ Crowley, \emph{Band-limitedness of the warped pullback \(Y\)}, this repository (\texttt{docs/proof\_work/BandLimitedness.tex}); full mode-sum estimates and the weak-\(*\) spectral-measure argument that complement Theorem~\ref{thm:bandlimit} above.
\bibitem{Gabcke} W.\ Gabcke, \emph{Neue Herleitung und explizite Restabsch\"atzung der Riemann--Siegel-Formel}, Dissertation, G\"ottingen, 1979 (reissued 2015); full asymptotic expansion of \(R(\tau)\) in \(\tau^{-1/2}\) with \(1\)-periodic real-analytic coefficients.
\bibitem{Berry} M.V.\ Berry, \emph{The Riemann--Siegel expansion for the zeta function: high orders and remainders}, Proc.\ R.\ Soc.\ Lond.\ A \textbf{450} (1995), 439--462.
\bibitem{BerryKeating} M.V.\ Berry and J.P.\ Keating, \emph{A new asymptotic representation for \(\zeta(\tfrac12+it)\) and quantum spectral determinants}, Proc.\ R.\ Soc.\ Lond.\ A \textbf{437} (1992), 151--173.
\bibitem{Edwards} H.M.\ Edwards, \emph{Riemann's Zeta Function}, Academic Press, 1974, \S7.4 (Riemann--Siegel contour-integral representation).
\bibitem{Hormander} L.\ H\"ormander, \emph{The Analysis of Linear Partial Differential Operators I}, Springer, 1983, Theorem 7.3.1 (Schwartz--Paley--Wiener).
\bibitem{Stein} E.M.\ Stein, \emph{Harmonic Analysis: Real-Variable Methods, Orthogonality, and Oscillatory Integrals}, Princeton, 1993, Chap.~VIII \S1 (non-stationary phase).
\end{thebibliography}

\end{document}
